\documentclass[11pt,a4paper]{scrartcl}

% -------------------------------------------------------
% Pakete
% -------------------------------------------------------
\usepackage[utf8]{inputenc}
\usepackage[T1]{fontenc}
\usepackage[ngerman]{babel}
\usepackage{lmodern}
\usepackage[left=2cm, right=2cm, top=1.8cm, bottom=1.8cm]{geometry}
\usepackage{amsmath, amssymb}
\usepackage{xcolor}
\usepackage{enumitem}
\usepackage{multicol}
\usepackage{array}
\usepackage{booktabs}
\usepackage{listings}
\usepackage{needspace}

\pagestyle{empty}
\setlength{\parindent}{0pt}
\setlength{\parskip}{3pt}

% Punkte am rechten Rand
\newcommand{\punkte}[1]{\hfill\mbox{\small(#1\,P)}}

% -------------------------------------------------------
% Formularfelder (digitale Fassung)
% -------------------------------------------------------
% unicode=false: sonst bekommen Optionsfelder (radio) kaputte Namen (hyperref 7.01)
\usepackage[unicode=false, pdftitle={Leistungskontrolle Informatik}]{hyperref}
\definecolor{feldRand}{HTML}{005A9C}
% Aussehen aller Felder: blauer Rand, hellblauer Hintergrund, Schrift 11 pt
\renewcommand{\DefaultOptionsofText}{print, bordercolor={0 0.353 0.612},
  backgroundcolor={0.94 0.96 1}, charsize=11pt}
\renewcommand{\DefaultOptionsofCheckBox}{print, bordercolor={0 0.353 0.612},
  backgroundcolor={0.94 0.96 1}, charsize=11pt}
\renewcommand{\DefaultHeightofText}{0.6cm}
% Jedes Feld steckt in einer eigenen Box (raisebox/hbox). Sonst übernimmt pdfTeX die
% Höhe der ganzen Tabellenzeile und die Felder werden zu hoch.
% \feld[Breite]{Name}      einzeiliges Textfeld (statt Schreiblücke)
\newcommand{\feld}[2][2.4cm]{\raisebox{-0.17cm}{\TextField[name=#2, width=#1]{}}}
% \zeichen{Name}           Feld für genau ein Zeichen, zentriert
\newcommand{\zeichen}[1]{\raisebox{-0.12cm}{\TextField[name=#1, width=0.8cm, height=0.5cm, maxlen=1, align=1]{}}}
% \kreuz{Name}             Ankreuzkästchen (statt $\square$)
\newcommand{\kreuz}[1]{\raisebox{-0.08cm}{\CheckBox[name=#1, width=0.42cm, height=0.42cm]{}}}
% \antwort[Höhe]{Name}     mehrzeiliges Textfeld (statt Schreiblinien)
\newcommand{\antwort}[2][3.2cm]{\par\vspace{0.15cm}\noindent
  \hbox{\TextField[name=#2, multiline=true, width=\linewidth, height=#1]{}}}

% -------------------------------------------------------
% Python-Listings
% -------------------------------------------------------
\definecolor{pyBack}{HTML}{F7F8FA}
\definecolor{pyRule}{HTML}{005A9C}
\definecolor{pyKey}{HTML}{0B4F8A}
\definecolor{pyStr}{HTML}{A3361B}
\definecolor{pyCom}{HTML}{5F7A5F}
\lstdefinestyle{pythonMKG}{
  language=Python,
  basicstyle=\ttfamily\small,
  keywordstyle=\color{pyKey}\bfseries,
  stringstyle=\color{pyStr},
  commentstyle=\color{pyCom}\itshape,
  backgroundcolor=\color{pyBack},
  frame=leftline, framerule=2pt, rulecolor=\color{pyRule},
  xleftmargin=8pt, framexleftmargin=6pt,
  showstringspaces=false, columns=fullflexible, keepspaces=true, upquote=true,
  numbers=left, numberstyle=\tiny\color{gray}, numbersep=12pt,
  aboveskip=4pt, belowskip=4pt,
}
\lstset{style=pythonMKG}

\setlist[enumerate,1]{label=(\arabic*), leftmargin=*, itemsep=10pt}
\setlist[enumerate,2]{label=(\alph*), itemsep=4pt}

\begin{document}
\shorthandoff{"}
\begin{Form}

% -------------------------------------------------------
% KOPF
% -------------------------------------------------------
Name: \feld[6cm]{name} \hfill Klasse: \feld[2cm]{klasse} \hfill Datum: \feld[3cm]{datum}

\vspace{0.4cm}
\begin{center}
  {\Large\bfseries Leistungskontrolle Informatik}\\[0.15cm]
  {\large Python-Grundlagen: Programme, Daten und Datentypen}
\end{center}

{\small\itshape Klicke in die blauen Felder und schreibe oder kreuze dort deine Antworten an.
Speichere die Datei vor dem Hochladen und öffne sie zur Kontrolle noch einmal.}

% -------------------------------------------------------
\subsection*{A \quad Programme und Daten}
% -------------------------------------------------------

\begin{enumerate}

\needspace{4cm}
\item Kreuze an, welche Zeile des Programms zur Verarbeitung gehört. \punkte{1}
\begin{lstlisting}
seite = float(input("Seitenlaenge in m: "))
flaeche = seite * seite
print("Flaeche:", flaeche, "m2")
\end{lstlisting}
\kreuz{a1-z1}\ Zeile 1 \qquad \kreuz{a1-z2}\ Zeile 2 \qquad \kreuz{a1-z3}\ Zeile 3

\needspace{6cm}
\item Ein Getränkeautomat arbeitet nach dem EVA-Prinzip. Ordne jeden Schritt der Eingabe (E),
der Verarbeitung (V) oder der Ausgabe (A) zu. \punkte{3}

\renewcommand{\arraystretch}{1.35}
\begin{tabular}{c p{11.2cm} c}
\toprule
Schritt & Was passiert? & E / V / A \\
\midrule
1 & Das Display zeigt \glqq{}Bitte Getränk entnehmen\grqq{} an. & \zeichen{a2-s1} \\
2 & Du wirfst Münzen ein. & \zeichen{a2-s2} \\
3 & Der Automat berechnet das Wechselgeld. & \zeichen{a2-s3} \\
4 & Du wählst per Tastendruck ein Getränk aus. & \zeichen{a2-s4} \\
5 & Die Flasche fällt in das Ausgabefach. & \zeichen{a2-s5} \\
6 & Der Automat zählt den Wert der eingeworfenen Münzen zusammen. & \zeichen{a2-s6} \\
\bottomrule
\end{tabular}

\needspace{6cm}
\item Die Zeilen eines Programms sind durcheinandergeraten. Das Programm soll zuerst
nach dem Widerstand fragen. Gib die richtige Reihenfolge der Buchstaben an. \punkte{2}

\begin{tabular}{>{\bfseries}l >{\ttfamily}l}
A & u = r * i \\
B & print("Spannung:", u, "V") \\
C & \# Spannung nach dem Ohmschen Gesetz \\
D & i = float(input("Stromstaerke in A: ")) \\
E & r = float(input("Widerstand in Ohm: ")) \\
\end{tabular}

\vspace{0.2cm}
Reihenfolge: \ \zeichen{a3-1}\ \zeichen{a3-2}\ \zeichen{a3-3}\ \zeichen{a3-4}\ \zeichen{a3-5}

\needspace{5cm}
\item Entscheide für jede Aussage, ob sie richtig oder falsch ist. \punkte{2}

\renewcommand{\arraystretch}{1.3}
\begin{tabular}{p{12.2cm}cc}
\toprule
Aussage & richtig & falsch \\
\midrule
Eine Zeile, die mit \texttt{\#} beginnt, ist ein Kommentar und wird beim Ausführen
übersprungen. & \kreuz{a4-1-r} & \kreuz{a4-1-f} \\
Nach \texttt{x = 5} und danach \texttt{x = 8} hat \texttt{x} den Wert 13. & \kreuz{a4-2-r} & \kreuz{a4-2-f} \\
Bei \texttt{a = b} wird der Wert von \texttt{b} in \texttt{a} kopiert, \texttt{b} bleibt
unverändert. & \kreuz{a4-3-r} & \kreuz{a4-3-f} \\
\texttt{input()} gibt einen Text auf dem Bildschirm aus, ohne auf eine Eingabe zu warten. &
\kreuz{a4-4-r} & \kreuz{a4-4-f} \\
\bottomrule
\end{tabular}

\needspace{5.5cm}
\item Nenne vier Regeln, die du beim Festlegen eines Variablennamens in Python beachten musst.
\punkte{4}
\antwort{a5}

\needspace{5cm}
\item Erkläre den Unterschied zwischen einem Syntaxfehler und einem Semantikfehler.
\punkte{2}
\antwort{a6}

\end{enumerate}

% -------------------------------------------------------
\needspace{9cm}
\subsection*{B \quad Datentypen und Operatoren}
% -------------------------------------------------------

\begin{enumerate}[resume]

\needspace{5.5cm}
\item Kreuze für jeden Wert den passenden Datentyp an. \punkte{3}

\renewcommand{\arraystretch}{1.25}
\begin{tabular}{>{\ttfamily}l ccc}
\toprule
\textrm{Wert} & \texttt{int} & \texttt{float} & \texttt{str} \\
\midrule
15 & \kreuz{a7-1-int} & \kreuz{a7-1-float} & \kreuz{a7-1-str} \\
1.72 & \kreuz{a7-2-int} & \kreuz{a7-2-float} & \kreuz{a7-2-str} \\
"04109" & \kreuz{a7-3-int} & \kreuz{a7-3-float} & \kreuz{a7-3-str} \\
2.0 & \kreuz{a7-4-int} & \kreuz{a7-4-float} & \kreuz{a7-4-str} \\
"7" & \kreuz{a7-5-int} & \kreuz{a7-5-float} & \kreuz{a7-5-str} \\
-3 & \kreuz{a7-6-int} & \kreuz{a7-6-float} & \kreuz{a7-6-str} \\
\bottomrule
\end{tabular}

\needspace{2.5cm}
\item Gib an, was Python bei den folgenden Rechnungen ausgibt. \punkte{3}

\begin{multicols}{3}
\begin{enumerate}
  \item \texttt{17 // 5} \ \feld[1.5cm]{a8-a}
  \item \texttt{17 \% 5} \ \feld[1.5cm]{a8-b}
  \item \texttt{10 / 4} \ \feld[1.5cm]{a8-c}
\end{enumerate}
\end{multicols}

\end{enumerate}

% -------------------------------------------------------
\needspace{6cm}
\subsection*{Zusatzaufgabe}
% -------------------------------------------------------

Mia und Ben sollen die Plätze tauschen. Ergänze die Lücken so, dass der Tausch mit einer
Hilfsvariable gelingt. \punkte{+3}

\begin{tabular}{>{\ttfamily}l}
platz1 = "Mia" \\
platz2 = "Ben" \\
hilf = \ \feld[3cm]{z-1} \\[0.25cm]
platz1 = \ \feld[3cm]{z-2} \\[0.25cm]
platz2 = \ \feld[3cm]{z-3} \\
\end{tabular}

\vspace{0.6cm}
\begin{center}
\small Viel Erfolg!
\end{center}

\end{Form}
\end{document}
